\phantomsection
\addcontentsline{toc}{subsection}{Fundamentos Previos}
\fontsize{12}{18}\selectfont{\textcolor{blueunam}{\textbf{QUÍMICA DE MATERIALES}}}
\subsection*{\textcolor{blueunam}{Fundamentos Previos}}

\phantomsection
\addcontentsline{toc}{subsubsection}{Tipos de enlace}
\subsubsection*{\textcolor{blueunam}{Tipos de enlace}}

\textbf{Enlace covalente:} interacción entre dos átomos en la que sus 
orbitales atómicos se solapan en fase, sumando constructivamente las 
funciones de onda y aumentando la densidad electrónica entre los núcleos. 
Esa densidad es atraída por ambos núcleos positivos, bajando la energía 
potencial del sistema. Si esa estabilización supera la repulsión 
núcleo-núcleo, el sistema se estabiliza y existe el enlace. Los átomos 
comparten uno o más pares de electrones, generalmente uno por átomo. 
Se considera covalente cuando la diferencia de electronegatividad es 
menor a 1.7.

\begin{center}
\schemestart
\chemname{\chemfig{H-C(-[2]H)(-[6]H)-H}}{Metano \ce{CH4}}
\schemestop
\end{center}

\textbf{Enlace sencillo:} enlace covalente en el que dos átomos comparten 
exactamente un par de electrones mediante solapamiento frontal de orbitales, 
formando un enlace $\sigma$. Es el enlace de menor orden (O.E.$=1$) y 
permite rotación libre alrededor del eje internuclear.

\begin{center}
\schemestart
\chemname{\chemfig{H-H}}{\ce{H2}, enlace $\sigma$}
\schemestop
\end{center}

\textbf{Enlace polar:} enlace covalente en el que los electrones no se 
distribuyen simétricamente porque los átomos tienen distinta 
electronegatividad. El más electronegativo concentra mayor densidad 
electrónica ($\delta^-$) y el otro queda con una deficiencia ($\delta^+$), 
generando un dipolo de enlace. Si los dipolos se cancelan vectorialmente 
por la geometría molecular, la molécula tiene momento dipolar neto igual 
a cero.

\begin{center}
\schemestart
\chemname{\chemfig{H-[104.5]\charge{45=\:,135=\:}{O}-[-104.5]H}}{Ejemplo: \ce{H2O}, geometría angular, dipolos no se cancelan, $\mu \neq 0$}
\schemestop
\end{center}

\begin{center}
\schemestart
\chemname{\chemfig{O=C=O}}{Contraejemplo: \ce{CO2}, geometría lineal, dipolos se cancelan, $\mu = 0$}
\schemestop
\end{center}

\textbf{Estructura de Lewis:} representación que muestra explícitamente 
todos los electrones de valencia de una molécula o ion. Los compartidos 
se representan como líneas entre átomos, cada una valiendo 2 electrones, 
y los no compartidos como pares libres sobre cada átomo. Se construye 
distribuyendo los electrones de valencia totales minimizando cargas formales.

\begin{center}
\schemestart
\chemname{\chemfig{H-\charge{90=\:}{N}(-[6]H)-H}}{\ce{NH3}: $e^-_T = 5+(1\times3) = 8$, C.F. N $= 5-2-3 = 0$, C.F. H $= 0$}
\schemestop
\end{center}


\phantomsection
\addcontentsline{toc}{subsubsection}{Carga formal vs densidad electrónica real}
\subsubsection*{\textcolor{blueunam}{Carga formal vs densidad electrónica real}}

La \textbf{carga formal} es un conteo entero: cuántos electrones completos 
le sobran o faltan a un átomo respecto a su estado neutro, asumiendo que 
los electrones enlazantes se dividen exactamente a la mitad. Es una 
convención útil para verificar estructuras de Lewis, no una medida de 
distribución real.

En realidad los electrones no ``pertenecen'' fijamente a ningún átomo: 
son probabilísticos. Lo que sí es real es que estadísticamente, en promedio, 
ciertos átomos concentran más densidad electrónica que otros. Eso se 
describe con la carga parcial $\delta$: un átomo con más densidad promedio 
tiene $\delta^-$, uno con menos tiene $\delta^+$. Si la diferencia es mayor, 
la carga parcial es mayor, pero siempre es un promedio estadístico, nunca 
un valor discreto fijo como en la carga iónica real.

\textbf{En enlaces iónicos} como \ce{NaCl}, la carga sí es real: el 
\ce{Na+} físicamente perdió un electrón y tiene 10 electrones. En 
covalente, la carga formal \ce{+1} o \ce{-1} no implica transferencia 
real, solo un déficit en el conteo convencional.

\begin{center}
\small
\renewcommand{\arraystretch}{1.4}
\begin{tabular}{|L{3.5cm}|L{4cm}|L{4cm}|}
\hline
\textbf{Concepto} & \textbf{Qué describe} & \textbf{Ejemplo} \\
\hline
Carga formal & Conteo entero de electrones de más o de menos & 
\ce{O} en \ce{OH-}: C.F. $= -1$ \\
\hline
Carga parcial $\delta$ & Probabilidad de encontrar densidad en un átomo & 
\ce{C-O}: $\delta^+$ en C, $\delta^-$ en O \\
\hline
Carga iónica real & Transferencia física de un electrón completo & 
\ce{Na+} en \ce{NaCl}: realmente tiene 10 $e^-$ \\
\hline
\end{tabular}
\normalsize
\end{center}